\chapter{Discussion}
\label{ch:discussion}

This chapter interprets the findings of Chapter~4 against the research questions of Section~\ref{sec:aims} and the literature of Chapter~2, examines the limitations of the study, and draws out the methodological lessons with relevance beyond this specific system.

\section{Synthetic Supervision and Spatial Grounding (RQ1)}
\label{sec:disc_rq1}

The central question of this dissertation was whether cheaply composited synthetic data can teach a VLA to spatially ground transparent objects. The answer the results support is a qualified yes: qualified not by the model, but by the data. The same architecture, optimiser, and budget produced two opposite outcomes---complete shortcut memorisation in the defective configuration, and genuine image-conditioned prediction (image-conditioned, per-sample predictions tracking target positions, in contrast to the constant output of the defective configuration) once per-sample target diversity, background-identity splitting, and compositing-seam removal were enforced.

This finding gives empirical form to a hazard the synthetic-data literature under-documents for \emph{action} supervision (Section~\ref{sec:lit_synthetic}). When the supervised quantity is a class label, shortcut learning degrades accuracy visibly. When it is a spatial coordinate, shortcut learning can hide behind flattering aggregate metrics: the memorising policy scored 97.7\% token accuracy because five of seven action dimensions were constants, and its training loss reached exactly zero. Both signals would pass a superficial review. The diagnostic that exposed it---per-coordinate error on held-out data, plus treating zero training loss as an alarm rather than an achievement---is a transferable protocol for anyone fine-tuning VLAs on procedurally generated data, and complements benchmark-level observations that VLA spatial competence must be deliberately engineered rather than assumed~\citep{cai2024, yuan2024}.

The residual accuracy ceiling also merits interpretation. Both corrected training runs, despite differing learning rates and stopping schedules, converged to statistically indistinguishable validation losses and //// $\approx$28-bin spatial error////////. The bottleneck is therefore the data distribution, not optimisation: one object category, one instruction template, and 2D compositing without refraction or contact shadows bound what spatial precision the policy can extract. This is consistent with the direction of the spatial-grounding literature, which improves localisation through richer supervision or representations---dedicated pointing data~\citep{yuan2024}, explicit 3D encodings~\citep{qu2025, zhen2024}---rather than through longer training on the same data.

\section{Perception Generalisation and Error Propagation (RQ2, RQ3)}
\label{sec:disc_rq23}

The perception results reproduce, at smaller scale, the pattern that made ClearGrasp viable~\citep{sajjan2020}: RGB-only networks recover transparent-object geometry well within their training distribution, and their segmentation transfers to novel shapes far better than their normal estimation. The shape-dependence of the degradation---6$^{\circ}$ loss on a geometrically similar unseen object versus 14.5$^{\circ}$ on a dissimilar one---suggests that broadening the \emph{shape} distribution of training data matters more for geometric generalisation than the quantity of data on any single shape, a directly actionable conclusion for future data collection.

The ablation delivers the dissertation's cleanest quantitative contribution: the isolation of perception-to-depth error propagation. Because the solver with perfect normals statistically matches a strong inpainting baseline (median difference 3.5\,mm), the 0.28\,m RMSE deficit of the deployed configuration is attributable to the 19.55$^{\circ}$ normal error alone. This converts an abstract perception metric into an operational cost in metres, and identifies normal accuracy---not solver design---as the highest-leverage improvement target in the geometry path.

The finding that geometry-unaware inpainting is competitive on these objects deserves an honest reading rather than a defensive one. ClearGrasp's evaluation objects are shallow relative to their surroundings, so boundary interpolation recovers most of the depth field; the normal-guided machinery earns its complexity only where surfaces curve substantially. Reporting this boundary condition is itself a result: it delineates where the ClearGrasp-style pipeline is necessary and where a three-line baseline suffices---information a deployment engineer needs and the original literature does not emphasise.

 
\section{From Spatial Accuracy to Grasp Feasibility (RQ4)}
\label{sec:disc_rq4}
 
The grasp results tell a clear causal story across two training iterations.
The first-generation system failed at 2.6\% on-target, and three specific
data defects were identified: 10\% invalid training labels, missing
transparency cues in composited images, and destructive double
downsampling. The second iteration corrected all three and nearly tripled
the on-target rate to 7.2\%.
 
The type breakdown is the sharpest evidence in this dissertation. On native
ClearGrasp frames---which contain true refraction, cast shadows, and
specular highlights---the model achieves 11.3\% on-target and 20.2\%
on-any-object. On composites---which lack all three cues---it achieves only
3.1\% and 5.5\%. This four-to-one ratio isolates compositing fidelity as
the dominant cause of the spatial precision deficit, not the model
architecture, not the LoRA configuration, and not the training budget. The
composite performance in Run~2 is nearly identical to Run~1's overall
performance (2.6\%), indicating that the label and resolution fixes
contributed marginally compared to the addition of authentic visual cues.
 
The size breakdown provides a secondary diagnosis: 0\% on-target for objects
below 142 pixels of area, 14.1\% for objects above 562 pixels. At 224$\times$224,
small transparent objects are resolved by too few encoder patches to carry
localisation signal. This suggests that higher processing resolution would
yield further gains, at the cost of deviating from OpenVLA's pretrained
encoder configuration.
 
These findings have a practical implication for anyone considering
synthetic-only VLA training for transparent-object manipulation: compositing
is insufficient. The visual cues that distinguish transparent localisation
from opaque localisation---refraction distortion of the background, caustic
patterns, specular geometry---are exactly what compositing destroys.
Physically based rendering (as in the ClearGrasp assets used here) or real
captures are necessary training components, not optional enrichments. This
is the dissertation's most actionable engineering conclusion.
 
The 7.2\% overall on-target rate remains far below deployment requirements.
Whether the residual gap is closable within the OpenVLA architecture (with
higher-fidelity data, more categories, and higher resolution) or requires a
dedicated spatial-prediction architecture such as
RoboPoint~\citep{yuan2024} is the open question this work bequeaths to
future investigation.
 
The latency profile inverts a common assumption: the 7B language model
accounts for 35\% of the 3.87\,s frame time; the CPU-bound sparse depth
solver accounts for 65\%. Given the ablation finding that inpainting
outperforms the solver on shallow objects, replacing the solver with a
single \texttt{cv2.inpaint} call would halve the per-frame cost with no
loss of depth accuracy for the evaluated morphologies.
 
 
\section{Limitations}
\label{sec:disc_limitations}

The following limitations bound the claims of this work. (i)~All policy evaluation is on synthetic composites; no real-camera or physical-robot validation was performed, so the sim-to-real gap for the \emph{policy} is unquantified (the perception module inherits ClearGrasp's evidence of transfer, but this was not re-verified here). (ii)~Policy training covers a single object category and a single instruction template; claims about semantic flexibility remain untested pending the experiment of Section~\ref{sec:res_semantic}. (iii)~The in-distribution perception figures overlap the training data and are reported as upper bounds, with the OOD categories carrying the generalisation claim. (iv)~The grasp proxy is a geometric criterion, not a physics simulation: it ignores gripper geometry, object slip, and approach kinematics. (v)~Checkpoint-level policy evaluation used 50 samples per checkpoint for tractability; the headline grasp measurement used 197. (vi)~The action space fixes height and orientation, presuming top-down grasps on a planar belt; the pipeline does not address 6-DOF grasp synthesis.

\section{Summary}
\label{sec:disc_summary}

Taken together, the results characterise both the promise and the present limits of synthetic-only VLA training for transparent-waste sorting: the data pipeline, not the model, decides whether spatial learning happens at all; perception error is the dominant and now-quantified cost in the geometry path; and the distance to deployable grasp precision is measurable, explainable, and attributable to identifiable data-fidelity gaps rather than to any single component failure. The final chapter draws these threads into conclusions and recommendations.